\markdownRendererDocumentBegin
\markdownRendererSectionBegin
\markdownRendererSectionBegin
\markdownRendererHeadingTwo{\markdownRendererStrongEmphasis{1.6\markdownRendererCircumflex\markdownRendererCircumflex I The Contents and Purposes of this Handbook}}\markdownRendererInterblockSeparator
{}\markdownRendererUlBeginTight
\markdownRendererUlItem [\markdownRendererEmphasis{see more details in \markdownRendererLink{notes while reading reviewer comments}{https://docs.google.com/document/d/1PKhH5R25GfpMx1LtoxfF316u60uLDnYpyy8YTngGxVs/edit?tab=t.0}{https://docs.google.com/document/d/1PKhH5R25GfpMx1LtoxfF316u60uLDnYpyy8YTngGxVs/edit?tab=t.0}{}}]\markdownRendererUlItemEnd 
\markdownRendererUlItem MAKE SURE TO SIGNPOST THIS SECTION THROUGHOUT CH 1\markdownRendererUlItemEnd 
\markdownRendererUlEndTight \markdownRendererInterblockSeparator
{}\markdownRendererSectionBegin
\markdownRendererHeadingThree{\markdownRendererStrongEmphasis{GENERAL}}\markdownRendererInterblockSeparator
{}\markdownRendererUlBeginTight
\markdownRendererUlItem We should create a new section in Ch1, something like "The purpose, structure, and content of this handbook" that deals in many of the high-level, purpose-of-polar issues that come up from reviewers. \markdownRendererStrongEmphasis{A running list of these:}\markdownRendererInterblockSeparator
{}\markdownRendererUlBeginTight
\markdownRendererUlItem This is a tool, for medium granularity phonetic transcription\markdownRendererInterblockSeparator
{}\markdownRendererUlBeginTight
\markdownRendererUlItem PoLaR (like much of annotation) uses discrete symbols; despite this, phonetics is not phonological/symbolic — this is a common mistake in interpretation of any phonetic transcription (and not a problem of the practice of transcribing phonetics with discrete symbols)\markdownRendererUlItemEnd 
\markdownRendererUlItem Other common intonational annotation tools (e.g., the ToBI family) are both phonological analyses as well as annotated observations; but it is not necessary for intonational annotation systems to do both. PoLaR is not a phonological analysis, but rather its annotated observations can still INFORM the phonological analysis.\markdownRendererInterblockSeparator
{}\markdownRendererUlBeginTight
\markdownRendererUlItem FURTHERMORE, not all of the perceptible/relevant intonational modulations are phonological (cf 'emotional prosody')\markdownRendererUlItemEnd 
\markdownRendererUlItem we want to track observations of any type of intonational modulation\markdownRendererUlItemEnd 
\markdownRendererUlEndTight \markdownRendererUlItemEnd 
\markdownRendererUlItem Let's figure out where to say this---maybe early on---and how to word it carefully\markdownRendererInterblockSeparator
{}\markdownRendererUlBeginTight
\markdownRendererUlItem maybe relating it to the fact that different speech users may use different cues?\markdownRendererUlItemEnd 
\markdownRendererUlItem PoLaR will have different annotations when speakers use different phonetic cues for the same phonological category (tracking variation!)\markdownRendererUlItemEnd 
\markdownRendererUlItem because different listeners rely on different cues, you can get different annotations\markdownRendererUlItemEnd 
\markdownRendererUlEndTight \markdownRendererUlItemEnd 
\markdownRendererUlEndTight \markdownRendererUlItemEnd 
\markdownRendererUlItem General point about phonetic annotation (potentially interesting): \markdownRendererStrongEmphasis{human intuition is used in all levels of phonetic annotation} (how and whether to annotate) because phonetic annotation is just a registering of directly-observable values\markdownRendererInterblockSeparator
{}\markdownRendererUlBeginTight
\markdownRendererUlItem so the purpose of this handbook is not intending to provide  “the” set of labels for any example\markdownRendererInterblockSeparator
{}\markdownRendererUlBeginTight
\markdownRendererUlItem they can’t just write ANYTHING; differences are limited in possibilities by this handbook\markdownRendererInterblockSeparator
{}\markdownRendererUlBeginTight
\markdownRendererUlItem there are wrong labels; disagreements are acceptable but there are still mistakes\markdownRendererUlItemEnd 
\markdownRendererUlEndTight \markdownRendererUlItemEnd 
\markdownRendererUlItem this handbook is here to ensure that when labellers differ, it’s because they PERCEIVE differences\markdownRendererInterblockSeparator
{}\markdownRendererUlBeginTight
\markdownRendererUlItem as with any annotation system of this sort, it’s probably the case that novices will be divergent at first, but through practice and review of these guidelines, they should converge closer together (though should not be expected to end up with 100\markdownRendererPercentSign{} —or even a majority of— identical labels)\markdownRendererUlItemEnd 
\markdownRendererUlItem there are some contours where it will be easier to get alignment between annotators (e.g., a straight line, vs one with lots of slow curves)\markdownRendererUlItemEnd 
\markdownRendererUlEndTight \markdownRendererUlItemEnd 
\markdownRendererUlEndTight \markdownRendererUlItemEnd 
\markdownRendererUlItem PoLaR is for descriptive observations and NOT categorical/phonological analysis. we EXPECT people to make different observations for the same data, at different levels of granularity (similar to narrow transcription IPA)\markdownRendererInterblockSeparator
{}\markdownRendererUlBeginTight
\markdownRendererUlItem (about labeller intuitions, and how they might truly vary at linguistic/cognitive level — i.e., variation is not always noise)\markdownRendererStrongEmphasis{.}\markdownRendererUlItemEnd 
\markdownRendererUlEndTight \markdownRendererUlItemEnd 
\markdownRendererUlEndTight \markdownRendererUlItemEnd 
\markdownRendererUlItem we provide guidance on workflows to increase informativeness of labels, because IRR is not necessarily a goal\markdownRendererInterblockSeparator
{}\markdownRendererUlBeginTight
\markdownRendererUlItem \markdownRendererStrongEmphasis{We don’t want to buy replicability at the cost of loss of information}\markdownRendererUlItemEnd 
\markdownRendererUlItem Variation is not always noise\markdownRendererInterblockSeparator
{}\markdownRendererUlBeginTight
\markdownRendererUlItem here are two sets of labels from expert labellers\markdownRendererUlItemEnd 
\markdownRendererUlItem they differ, let’s talk about how could you deal with these divergent labels\markdownRendererInterblockSeparator
{}\markdownRendererUlBeginTight
\markdownRendererUlItem detail-reduction process (e.g., “binning” within-labeller differences under a single label: ?* and * and ** all reduced to *)\markdownRendererUlItemEnd 
\markdownRendererUlItem consensus-building practices (e.g., with the PoLaR plugin)\markdownRendererUlItemEnd 
\markdownRendererUlEndTight \markdownRendererUlItemEnd 
\markdownRendererUlEndTight \markdownRendererUlItemEnd 
\markdownRendererUlItem \markdownRendererStrongEmphasis{hearken back to this point early on in 2.1.2 Points Tier}\markdownRendererUlItemEnd 
\markdownRendererUlEndTight \markdownRendererUlItemEnd 
\markdownRendererUlItem Talking-points about how to interpret (or not) PoLaR labels\markdownRendererInterblockSeparator
{}\markdownRendererUlBeginTight
\markdownRendererUlItem As a reader / analyst of PoLaR labels, it is worth being open to the possibility that the labeller (a) got it wrong \markdownRendererEmphasis{\markdownRendererStrongEmphasis{or}} (b) has a different intuition about the speech. So this means that interpreting phonetic labels can be a challenge (especially because of the role of intuition in doing labelling AND because of interspeaker variation in categories and associations between cues+categories)\markdownRendererUlItemEnd 
\markdownRendererUlItem Be especially cautious with phrase-final annotations, especially on the Points tier, and especially when the f0 track is unreliable\markdownRendererInterblockSeparator
{}\markdownRendererUlBeginTight
\markdownRendererUlItem Labellers should be familiar with pitch tracking errors; see section 2.2.2 (and review example files with this characteristic)\markdownRendererUlItemEnd 
\markdownRendererUlEndTight \markdownRendererUlItemEnd 
\markdownRendererUlEndTight \markdownRendererUlItemEnd 
\markdownRendererUlEndTight \markdownRendererUlItemEnd 
\markdownRendererUlEndTight \markdownRendererInterblockSeparator
{}
\markdownRendererSectionEnd \markdownRendererSectionBegin
\markdownRendererHeadingThree{\markdownRendererStrongEmphasis{RANGES}}\markdownRendererInterblockSeparator
{}\markdownRendererUlBeginTight
\markdownRendererUlItem My first and main issue regards the Pitch Ranges tier. To summarize, a) on the x-axis (temporal dimension), labellers identify the scope of such a pitch range for different portions of an utterance. Importantly, the scope of a given pitch range is independent of phrase breaks and may be assigned to any part of an utterance. Also, b) on the y-axis, annotators identify the f0 range (min / max) of a portion of an utterance and label the min and max on a tier. Min and max do not necessary refer to the actual minimum and maximum f0 value observed in the signal, but are values that the labellers determine as minimum and maximum for a particular speaker. Regarding a), allowing misaligned f0 ranges and phrases is difficult both from a theoretical and practical point of view and – in my opinion – not sufficiently motivated in the manuscript.\markdownRendererUlItemEnd 
\markdownRendererUlItem "but are values that the labellers determine as minimum and maximum for a particular speaker"\markdownRendererInterblockSeparator
{}\markdownRendererUlBeginTight
\markdownRendererUlItem we should focus on the fact that —and loudly broadcast it— that we're not trying to prove any position on a theory — instead, we are arguing that it's necessary or important to label these things, whose systematicity would be opaque if we weren't tracking them: and PoLaR is the tool for that — and those annotations could then be used as part of an argument for a position on a theory, but we're NOT doing that here\markdownRendererUlItemEnd 
\markdownRendererUlItem we’re trying to provide a tool to assess truth of claims about the relationship between acoustic signal and phonology; we are NOT making any particular claims about English or its prosodic phonology (or any prosodic phonology)\markdownRendererInterblockSeparator
{}\markdownRendererUlBeginTight
\markdownRendererUlItem though all transcription systems rely on some amount of analysis, we are trying to distance PoLaR from the specifics of any prosodic phonology\markdownRendererUlItemEnd 
\markdownRendererUlItem (but of course we do appeal to some of it; we build on the theoretical constructs of pitch accents and phrase boundaries)\markdownRendererUlItemEnd 
\markdownRendererUlEndTight \markdownRendererUlItemEnd 
\markdownRendererUlEndTight \markdownRendererUlItemEnd 
\markdownRendererUlItem "allowing misaligned f0 ranges and phrases is difficult both from a theoretical and practical point of view and – in my opinion – not sufficiently motivated in the manuscript."\markdownRendererInterblockSeparator
{}\markdownRendererUlBeginTight
\markdownRendererUlItem it's not motivated [by theory]! it's a proposal of what to track!\markdownRendererUlItemEnd 
\markdownRendererUlItem do exploratory annotation systems need to be "motivated"?\markdownRendererInterblockSeparator
{}\markdownRendererUlBeginTight
\markdownRendererUlItem we don't think so, but we \markdownRendererStrongEmphasis{should acknowledge that we realize that it's unprecedented in labelling}\markdownRendererUlItemEnd 
\markdownRendererUlItem it's also certainly not entirely off-the-wall, but it IS important to include in this system \markdownRendererStrongEmphasis{because of the data we have seen}.\markdownRendererUlItemEnd 
\markdownRendererUlItem things to bring up as motivation for tracking Ranges:\markdownRendererInterblockSeparator
{}\markdownRendererUlBeginTight
\markdownRendererUlItem some within-phrase range changes are attested:\markdownRendererInterblockSeparator
{}\markdownRendererUlBeginTight
\markdownRendererUlItem pitch-accent-related downstep. (though this is a bigger step forward in what's possible to annotate / expect.)\markdownRendererUlItemEnd 
\markdownRendererUlItem upstep in English (currently not annotatable in MAE\markdownRendererUnderscore{}ToBI; also a recognized part of the phonology in some Spanish varieties)\markdownRendererUlItemEnd 
\markdownRendererUlEndTight \markdownRendererUlItemEnd 
\markdownRendererUlItem prosodic events that are typically associated with the end of a phrase (e.g., pauses + lengthening) are well attested within phrases as well (e.g., interruptions)\markdownRendererUlItemEnd 
\markdownRendererUlItem what we've observed in the data\markdownRendererInterblockSeparator
{}\markdownRendererUlBeginTight
\markdownRendererUlItem \markdownRendererStrongEmphasis{the wider range of examples we can add the better}\markdownRendererUlItemEnd 
\markdownRendererUlEndTight \markdownRendererUlItemEnd 
\markdownRendererUlEndTight \markdownRendererUlItemEnd 
\markdownRendererUlItem also: how will we know if we don't track ranges?\markdownRendererInterblockSeparator
{}\markdownRendererUlBeginTight
\markdownRendererUlItem one of the very useful functions of polar will allow for tests of unspoken/uninvestigated hunches:\markdownRendererUlItemEnd 
\markdownRendererUlEndTight \markdownRendererUlItemEnd 
\markdownRendererUlEndTight \markdownRendererUlItemEnd 
\markdownRendererUlItem what are the criteria to determine that the ranges have changed?\markdownRendererInterblockSeparator
{}\markdownRendererUlBeginTight
\markdownRendererUlItem for now, we are vague on that\markdownRendererInterblockSeparator
{}\markdownRendererUlBeginTight
\markdownRendererUlItem asking listeners to go with their gut / track whether Levels annotations match their intuitions\markdownRendererUlItemEnd 
\markdownRendererUlItem Similar to divergent narrow IPA transcriptions: both could be right in some way, or both could be wrong in some way, either way \markdownRendererEmphasis{something} is going on that deserves attention\markdownRendererUlItemEnd 
\markdownRendererUlItem divergence and convergence are both informative! both beckon researchers/analysts to look at the signal again\markdownRendererUlItemEnd 
\markdownRendererUlEndTight \markdownRendererUlItemEnd 
\markdownRendererUlItem eventually we want better guidance\markdownRendererInterblockSeparator
{}\markdownRendererUlBeginTight
\markdownRendererUlItem multiple labellers separately label and then discuss\markdownRendererUlItemEnd 
\markdownRendererUlEndTight \markdownRendererUlItemEnd 
\markdownRendererUlEndTight \markdownRendererUlItemEnd 
\markdownRendererUlItem big point:\markdownRendererInterblockSeparator
{}\markdownRendererUlBeginTight
\markdownRendererUlItem we're not looking for a system that will have labellers always produce the same labellers\markdownRendererUlItemEnd 
\markdownRendererUlItem this is VALUABLE because divergences are places we need to look more closely\markdownRendererUlItemEnd 
\markdownRendererUlEndTight \markdownRendererUlItemEnd 
\markdownRendererUlItem Ranges\markdownRendererInterblockSeparator
{}\markdownRendererUlBeginTight
\markdownRendererUlItem start at a [ and end at a ]\markdownRendererInterblockSeparator
{}\markdownRendererUlBeginTight
\markdownRendererUlItem well accepted by all; very well attested; can be captured by any intonational annotation system\markdownRendererUlItemEnd 
\markdownRendererUlEndTight \markdownRendererUlItemEnd 
\markdownRendererUlItem start at a [ and change mid-phrase\markdownRendererInterblockSeparator
{}\markdownRendererUlBeginTight
\markdownRendererUlItem e.g., downstep — everyone knows about this, but may not be thinking of it as a “range” change\markdownRendererUlItemEnd 
\markdownRendererUlItem acknowledged by some as a “compression of ranges” (also as a drop in the pitch ceiling) —***\markdownRendererNbsp{}add some citations here***\markdownRendererUlItemEnd 
\markdownRendererUlEndTight \markdownRendererUlItemEnd 
\markdownRendererUlItem start mid-phrase end in the middle of another phrase\markdownRendererInterblockSeparator
{}\markdownRendererUlBeginTight
\markdownRendererUlItem range changes fully decoupled from phrases\markdownRendererUlItemEnd 
\markdownRendererUlItem have not been studied — it is an open question as to how common these are and what sorts of analysis might be appropriate\markdownRendererUlItemEnd 
\markdownRendererUlEndTight \markdownRendererUlItemEnd 
\markdownRendererUlEndTight \markdownRendererUlItemEnd 
\markdownRendererUlEndTight \markdownRendererUlItemEnd 
\markdownRendererUlItem "from a theoretical point of view, I am not convinced how – within one and the same phrase – speakers may employ different ranges. The authors themselves argue that phrase breaks may trigger a change in the f0 range (p. 106); I am not sure which kind of processes would trigger a change in f0 range within one and the same word."\markdownRendererInterblockSeparator
{}\markdownRendererUlBeginTight
\markdownRendererUlItem Phonological\markdownRendererInterblockSeparator
{}\markdownRendererUlBeginTight
\markdownRendererUlItem pitch accents! H+!H*\markdownRendererUlItemEnd 
\markdownRendererUlItem other phonological objects? (planet Vulcan + newtonian orbit?)\markdownRendererUlItemEnd 
\markdownRendererUlItem (also, w/in word phrase-boundaries: faaaaan H* L-H\markdownRendererPercentSign{} tastic H* L-L\markdownRendererPercentSign{})\markdownRendererUlItemEnd 
\markdownRendererUlItem we don't know what else, but in order to know what would trigger them (or even if they occur!), we would need to transcribe range changes with the possibility of changing ranges anywhere. i.e., to have a phonological analysis, we would need annotated data.\markdownRendererUlItemEnd 
\markdownRendererUlItem these may be rare, even vanishingly, but that doesn’t mean they shouldn’t be analyzable\markdownRendererInterblockSeparator
{}\markdownRendererUlBeginTight
\markdownRendererUlItem cf. periodic table and non-existant-in-nature elements\markdownRendererUlItemEnd 
\markdownRendererUlItem also, exceptions are worthy of study because they help refine our understanding of the system\markdownRendererUlItemEnd 
\markdownRendererUlEndTight \markdownRendererUlItemEnd 
\markdownRendererUlEndTight \markdownRendererUlItemEnd 
\markdownRendererUlItem Disfluencies\markdownRendererInterblockSeparator
{}\markdownRendererUlBeginTight
\markdownRendererUlItem accidental: deserve study and tracking\markdownRendererUlItemEnd 
\markdownRendererUlItem planned (i.e., planned non-canonical mapping between text and prosody):\markdownRendererSoftLineBreak
{}FAN. TAS. TIC.\markdownRendererUlItemEnd 
\markdownRendererUlEndTight \markdownRendererUlItemEnd 
\markdownRendererUlEndTight \markdownRendererUlItemEnd 
\markdownRendererUlItem From a methodological perspective, this assumption or possibility in PoLaR might lead to inconsistencies in labelling. That is, I am concerned about both the validity of the range annotation and its reliability. In fact, even though I am quite experienced in annotation, I would not have placed all f0 ranges as done by the authors, e.g., on p. 46, 106, 108, 109. I thus wonder how reliable the positioning of these range spaces is, which in turn, has implications for the pitch levels, as they are informed by the pitch ranges.\markdownRendererUlItemEnd 
\markdownRendererUlItem Theory / Labelling\markdownRendererInterblockSeparator
{}\markdownRendererUlBeginTight
\markdownRendererUlItem It’s very clear that range changes can matter, linguistically.\markdownRendererUlItemEnd 
\markdownRendererUlItem all of the types of labels that PoLaR was designed to employ are motivated by questions that have arisen in intonational research\markdownRendererInterblockSeparator
{}\markdownRendererUlBeginTight
\markdownRendererUlItem Current labelling systems cannot capture these changes. But because PoLaR explicitly annotates potential range changes, it makes explicit the challenge in identifying where ranges change — which can lead to new discoveries.\markdownRendererUlItemEnd 
\markdownRendererUlItem We\markdownRendererUlItemEnd 
\markdownRendererUlEndTight \markdownRendererUlItemEnd 
\markdownRendererUlItem We acknowledge that this is difficult; and so does the field\markdownRendererInterblockSeparator
{}\markdownRendererUlBeginTight
\markdownRendererUlItem identifying where ranges have changed “instantaneously”/phonologically (downstep) is relatively straightforward; but where they are drifting downwards continuously/phonetically, it would be surprising if two labellers identified the same point in time to mark a shifted range\markdownRendererUlItemEnd 
\markdownRendererUlItem distinguishing downstep from downdrift is well-known as challenging\markdownRendererInterblockSeparator
{}\markdownRendererUlBeginTight
\markdownRendererUlItem and though this labelling may be challenging for these reasons, PoLaR provides a tool for systematically labelling range changes and those labels (alongside other labels / measures) can be used to further elucidate answers to the questions of how they differ\markdownRendererUlItemEnd 
\markdownRendererUlEndTight \markdownRendererUlItemEnd 
\markdownRendererUlEndTight \markdownRendererUlItemEnd 
\markdownRendererUlEndTight \markdownRendererUlItemEnd 
\markdownRendererUlItem Practically Speaking\markdownRendererInterblockSeparator
{}\markdownRendererUlBeginTight
\markdownRendererUlItem We expect inconsistencies and that’s okay because of how we expect labellers to:\markdownRendererInterblockSeparator
{}\markdownRendererUlBeginTight
\markdownRendererUlItem \markdownRendererStrongEmphasis{Approach ambiguous signals with caution and an open mind}\markdownRendererSoftLineBreak
{}There may be cases where labellers are unsure of whether something is a range change or not.\markdownRendererSoftLineBreak
{}Labellers can use heuristics such as "be careful about proposing range changes when not at a phrase boundary — it's not impossible, but check if it's necessary. (If you are familiar with downstep pitch accents, this is obviously a place where you would want within-phrase range changes.)".\markdownRendererUlItemEnd 
\markdownRendererUlItem \markdownRendererStrongEmphasis{Disagree and then have discussions over disagreements}\markdownRendererInterblockSeparator
{}\markdownRendererUlBeginTight
\markdownRendererUlItem Inconsistencies in labelling are expected in first-round annotation, and we encourage labellers to meet after their first round of annotation and find disagreements between their labels and take them as launching points for discussion / exploration / further analysis.\markdownRendererInterblockSeparator
{}\markdownRendererUlBeginTight
\markdownRendererUlItem [To add: discussion on how to build consensus; possibly alluding to training videos that we will build; mention plugin — including notes about “at this current time” —\markdownRendererNbsp{}maybe in an appendix]\markdownRendererUlItemEnd 
\markdownRendererUlEndTight \markdownRendererUlItemEnd 
\markdownRendererUlItem in addition: disagreements are to be expected because different labellers use/perceive prosody differently. (also, there's a background assumption that annotators ought to be tapping into the same linguistic structures [phonetic, phonological, etc.], but that's not obviously true! in fact, in many, it's obviously not true. and so if we are tracking details at the phonetic level, we should certainly expect differences in labels to be in line with annotators' linguistic intuitions.)\markdownRendererInterblockSeparator
{}\markdownRendererUlBeginTight
\markdownRendererUlItem a big point here: there is no 1 English intonation. no 1 Mainstream American English intonation. there is so much variation; we want to build a system that captures any possibly of language/dialect/register.\markdownRendererUlItemEnd 
\markdownRendererUlEndTight \markdownRendererUlItemEnd 
\markdownRendererUlEndTight \markdownRendererUlItemEnd 
\markdownRendererUlEndTight \markdownRendererUlItemEnd 
\markdownRendererUlEndTight \markdownRendererUlItemEnd 
\markdownRendererUlItem With respect to b) (y-axis), I also have a methodological and a theoretical concern: Methodologically, I think that inferring speakers’ actual highs and lows requires intensive experience and such inferences are particularly prone for interrater variability. More importantly, assigning pitch levels based on the min/max of the f0 range represents a normalization that gets rid of the actual f0 values – which in some cases might be relevant indeed. I totally see that in some cases, e.g., the cases of range compression (p. 104) it is straightforward to apply such a normalization, but there might be cases, e.g., in different registers such as infant-directed speech, in which information of the actual values of pitch targets is important (e.g., when voices go into falsetto).\markdownRendererUlItemEnd 
\markdownRendererUlItem we think this means we need to do a better job at communicating the motivation and uses-cases/purpose of the ranges tier.\markdownRendererInterblockSeparator
{}\markdownRendererUlBeginTight
\markdownRendererUlItem 1. Motivation:\markdownRendererInterblockSeparator
{}\markdownRendererUlBeginTight
\markdownRendererUlItem 1 and the same Hz value can be high or low, so having a way of tracking relative f0 value is critical\markdownRendererUlItemEnd 
\markdownRendererUlItem moreover, that it ought to be relative with respect to a locally-defined domain.\markdownRendererUlItemEnd 
\markdownRendererUlItem note: this is not to replace direct f0 measures, but rather to supplement.\markdownRendererUlItemEnd 
\markdownRendererUlEndTight \markdownRendererUlItemEnd 
\markdownRendererUlItem 2. Goals:\markdownRendererInterblockSeparator
{}\markdownRendererUlBeginTight
\markdownRendererUlItem in various approaches to speech data, it is common to analyze and report pitch range size and extremes; but as f0 measures do not always instantiate the min/max (e.g. H* H-H\markdownRendererPercentSign{}), we want a systematic way to encode intuitions. these annotations can then be used to report more specific and fine-tuned values for f0 range size / extremes.\markdownRendererUlItemEnd 
\markdownRendererUlItem ranges help normalize f0 values into a scale that is coarser than direct measures, but finer-grained that high/low: Levels labels, which divide the range space into 5 bands. these Levels labels are useful fo describing melodies and discussing f0 at a medium-granularity.\markdownRendererInterblockSeparator
{}\markdownRendererUlBeginTight
\markdownRendererUlItem \markdownRendererStrongEmphasis{granularity between f0 and phonological labels is something worth saying loud and often}\markdownRendererUlItemEnd 
\markdownRendererUlItem not *narrow* phonetic transcription of f0 (something closer to direct f0 measures) but also not phonological transcription; broad phonetic transcription\markdownRendererUlItemEnd 
\markdownRendererUlItem (and we expect differences in where people threshold broad vs narrow vs something in between)\markdownRendererUlItemEnd 
\markdownRendererUlEndTight \markdownRendererUlItemEnd 
\markdownRendererUlEndTight \markdownRendererUlItemEnd 
\markdownRendererUlEndTight \markdownRendererUlItemEnd 
\markdownRendererUlItem the concerns, as we understand them, center around inter-rater differences: ranges invite fine-grained annotations that would easily vary across labellers. but do these concerns merit changes, given the motivations/goals?\markdownRendererInterblockSeparator
{}\markdownRendererUlBeginTight
\markdownRendererUlItem \markdownRendererStrongEmphasis{the motivation and goals do not include getting agreement on exact labels.} we agree that we can't expect identical ranges labels across labelers, or even within-labeler across times.\markdownRendererInterblockSeparator
{}\markdownRendererUlBeginTight
\markdownRendererUlItem many people may have discomfort with labelling in a way others might not. to address this, you can think of PoLaR as a log of \markdownRendererStrongEmphasis{the annotator’s perception} (and not as some indisputable objective truth about the intonation) —\markdownRendererNbsp{}even when it’s ambiguous.\markdownRendererUlItemEnd 
\markdownRendererUlItem we’re missing things by virtue of not being willing to address and resolve disagreements\markdownRendererUlItemEnd 
\markdownRendererUlEndTight \markdownRendererUlItemEnd 
\markdownRendererUlItem instead, the target is to get labels that yield both \markdownRendererStrongEmphasis{similar-enough} measures of f0 size/extrema and identical Levels labels.\markdownRendererInterblockSeparator
{}\markdownRendererUlBeginTight
\markdownRendererUlItem since f0 size/extrema measures are inherently imprecise, we don't view variation as a problem. (people don’t always recognize that a single number from a computer with multiple decimals is still imprecise.)\markdownRendererUlItemEnd 
\markdownRendererUlEndTight \markdownRendererUlItemEnd 
\markdownRendererUlItem Since the amount of precision required to get the same Levels is rather low (since each level occupies 20\markdownRendererPercentSign{} of the size of the pitch range), and since labellers are encouraged to discuss with each other (and could use same-Levels as a metric for Ranges-label agreement), we don’t think this requires rethinking as much as it requires clarification on \markdownRendererStrongEmphasis{what} we’re thinking.\markdownRendererUlItemEnd 
\markdownRendererUlEndTight \markdownRendererUlItemEnd 
\markdownRendererUlItem Taken together, the Ranges Tier in my view allows for too many “degrees of freedom” (Roettger, 2019) for labellers. This, in turn, may decrease reproducibility and reliability, which brings me to my second point.\markdownRendererUlItemEnd 
\markdownRendererUlItem We don’t think the “reproducibility and reliability” critiques (which we agree with) require new arguments here, but rather a kind of summary\markdownRendererInterblockSeparator
{}\markdownRendererUlBeginTight
\markdownRendererUlItem In our experience and prediction, there’s enough replicability to be useful\markdownRendererUlItemEnd 
\markdownRendererUlItem Also, the ways in which labellers differ can be USEFUL and is worth putting down\markdownRendererUlItemEnd 
\markdownRendererUlItem This is the kind of tool where, like RPT, multiple labellers’ work is better\markdownRendererUlItemEnd 
\markdownRendererUlEndTight \markdownRendererUlItemEnd 
\markdownRendererUlEndTight 
\markdownRendererSectionEnd 
\markdownRendererSectionEnd 
\markdownRendererSectionEnd \markdownRendererDocumentEnd